\documentclass[12pt,a4paper]{article}
\usepackage{amsmath,amssymb,amsthm}
\usepackage{hyperref}
\usepackage{geometry}
\geometry{margin=2.5cm}
\usepackage{enumitem}

\newtheorem{theorem}{Theorem}[section]
\newtheorem{lemma}[theorem]{Lemma}
\newtheorem{corollary}[theorem]{Corollary}
\newtheorem{proposition}[theorem]{Proposition}
\theoremstyle{definition}
\newtheorem{definition}[theorem]{Definition}
\theoremstyle{remark}
\newtheorem{remark}[theorem]{Remark}

\title{Quantum Foundations from Recognition Science:\\
The Born Rule, Complex Hilbert Space, Entanglement,\\
and the Measurement Problem}

\author{Jonathan Washburn\\
Recognition Science Research Institute\\
Austin, Texas\\
\texttt{washburn.jonathan@gmail.com}}

\date{March 2026}

\begin{document}
\maketitle

\begin{abstract}
We derive the foundational structure of quantum mechanics from
Recognition Science (RS).  Five structural features, usually
postulated, are shown to follow from the cost functional
$J(x) = \tfrac{1}{2}(x + x^{-1}) - 1$ and the 8-tick epoch:
(1)~the Born rule $P = |\psi|^2$ from $J$-cost phase invariance;
(2)~complex Hilbert space from the 8-tick half-period
$e^{i\pi} = -1$; (3)~entanglement as shared ledger constraint,
yielding the Tsirelson bound $2\sqrt{2}$;
(4)~measurement dissolved as the discrete ledger commit at epoch
boundaries; (5)~the canonical commutation relation
$[\hat{x}, \hat{p}] = i\hbar$ from double-entry non-commutativity.
All are derived with zero free parameters.
\end{abstract}

%% ============================================================
\section{Introduction}
\label{sec:intro}
%% ============================================================

Quantum mechanics has a structure that works extraordinarily well
but whose foundations remain disputed~\cite{Bell1987}.  Five
questions have resisted resolution since the 1930s:
\begin{enumerate}[nosep]
  \item Why the Born rule $P = |\psi|^2$?
  \item Why complex Hilbert space (not real or quaternionic)?
  \item What is measurement, ontologically?
  \item What is entanglement?
  \item Why $[\hat{x}, \hat{p}] = i\hbar$?
\end{enumerate}

Recognition Science derives all five from the single primitive
$J(x) = \tfrac{1}{2}(x + x^{-1}) - 1$.

%% ============================================================
\section{The Born Rule from Phase Invariance}
\label{sec:born}
%% ============================================================

\begin{definition}[Recognition Rate]
The recognition rate for a state with amplitude $A \in \mathbb{C}$
is $R(A) = \lim_{T \to \infty} N(A, T)/T$, where $N(A, T)$ is
the number of recognition events in $[0, T]$.
\end{definition}

\begin{theorem}[Born Rule]
\label{thm:born}
Let $R: \mathbb{C} \to [0, \infty)$ satisfy:
\begin{enumerate}[label=(\roman*),nosep]
  \item Phase invariance: $R(Ae^{i\theta}) = R(A)$ for all
  $\theta$.
  \item Normalization: $R(1) = 1$.
  \item Continuity: $R$ is continuous.
\end{enumerate}
Then $R(A) = |A|^2$.
\end{theorem}

\begin{proof}
By (i), $R(A) = f(|A|)$ for some function $f$.  Consider $n$
equal-amplitude branches with $|A| = 1/\sqrt{n}$.  The total
recognition rate is $n \cdot f(1/\sqrt{n})$.  By (ii) applied
to the full state, $n \cdot f(1/\sqrt{n}) = 1$, giving
$f(1/\sqrt{n}) = 1/n$ for all positive integers $n$.  Setting
$r = 1/\sqrt{n}$: $f(r) = r^2$ for $r \in \{1/\sqrt{n}
: n \in \mathbb{N}\}$.  By (iii), $f$ is continuous and the
values $1/\sqrt{n}$ are dense in $[0, 1]$, so $f(r) = r^2$
everywhere.  Therefore $R(A) = |A|^2$.
\end{proof}

\begin{corollary}[Interference]
For $|\Psi\rangle = c_1|1\rangle + c_2|2\rangle$:
$P = |c_1 + c_2|^2 = |c_1|^2 + |c_2|^2 +
2\,\mathrm{Re}(c_1^* c_2)$.
\end{corollary}

\begin{remark}
This derivation is analogous to Gleason's theorem~\cite{Gleason}
but uses RS axioms (phase invariance from $J$-cost) instead of
measure-theoretic assumptions on the lattice of projections.
\end{remark}

%% ============================================================
\section{Complex Hilbert Space from the 8-Tick Epoch}
\label{sec:hilbert}
%% ============================================================

\begin{theorem}[Complex Hilbert Space]
\label{thm:complex}
The quantum state space is a complex Hilbert space $\mathcal{H}$
over $\mathbb{C}$, not over $\mathbb{R}$ or $\mathbb{H}$.
\end{theorem}

\begin{proof}
The 8-tick epoch has period $2^D = 8$ (for $D = 3$).  The phase
structure traverses $e^{2\pi i k/8}$ for $k = 0, 1, \ldots, 7$.
At the half-period ($k = 4$):
\begin{equation}
  e^{i\pi} = -1.
\end{equation}
This requires a number field containing $\sqrt{-1}$, forcing
$\mathbb{C}$.
\begin{itemize}[nosep]
  \item $\mathbb{R}$: Cannot represent $e^{i\pi} = -1$ as a
  continuous rotation from $+1$ without extending to $\mathbb{C}$.
  \item $\mathbb{C}$: Exactly captures the 8-tick phase structure
  ($e^{2\pi i k/8}$, half-period $= e^{i\pi} = -1$).
  \item $\mathbb{H}$: Quaternions would require 16-tick
  periodicity (quaternionic phase $e^{i\pi}e^{j\pi} = -1 \cdot
  (-1) = 1$ at 16 ticks), inconsistent with the forced 8-tick
  structure.
\end{itemize}
\end{proof}

\begin{corollary}[Euler's Identity]
$e^{i\pi} + 1 = 0$ is a structural necessity of the 8-tick epoch,
not a mathematical coincidence.
\end{corollary}

\begin{remark}
Recent experiments have ruled out real-number quantum
mechanics~\cite{RealQM2022}, confirming that $\mathbb{C}$ is
physically required.  RS explains \emph{why}.
\end{remark}

%% ============================================================
\section{The Measurement Problem Dissolved}
\label{sec:measurement}
%% ============================================================

\begin{definition}[Ledger Commit]
At every 8-tick epoch boundary ($t \bmod 8 = 0$), the ledger
transitions from an uncommitted state (superposition of branches)
to a committed state (single definite outcome).  This is
``measurement.''
\end{definition}

\begin{theorem}[Measurement = Recognition]
\label{thm:measurement}
In RS, measurement is the discrete ledger commit.  No separate
measurement postulate is needed.
\end{theorem}

\begin{proof}
The uncommitted ledger state is
$|\Psi\rangle = \sum_i c_i |b_i\rangle$.  The committed state
is $|b_k\rangle$ with probability $|c_k|^2$ (by
Theorem~\ref{thm:born}).  The commit is irreversible: reversing
it would create an unbalanced ledger entry with $J \to \infty$.
Von Neumann's Process~1 (projection) and Process~2 (unitary
evolution) are unified: Process~2 is the inter-tick dynamics;
Process~1 is the epoch-boundary commit.
\end{proof}

\begin{proposition}[Definite Outcomes]
Each committed ledger entry is integer-valued.  No superposition
of committed entries exists.  Schr\"odinger's cat is never in a
superposition.
\end{proposition}

\begin{remark}
This resolves the Wigner's friend paradox: both observers read the
same committed ledger entry.  Wigner's pre-query ignorance is
epistemic, not ontic.
\end{remark}

%% ============================================================
\section{Entanglement as Shared Ledger Constraint}
\label{sec:entanglement}
%% ============================================================

\begin{definition}[Entanglement]
Two recognition events $A$ and $B$ are entangled if their joint
ledger state has a conserved quantity $Q_{AB} = Q_A + Q_B = q_0$
that constrains the individual states.
\end{definition}

\begin{theorem}[Bell Inequality Violation]
\label{thm:bell}
The shared ledger constraint produces correlations that violate
Bell's inequality.  The maximum violation is the Tsirelson bound:
\begin{equation}
  |\langle AB \rangle + \langle AB' \rangle +
  \langle A'B \rangle - \langle A'B' \rangle|
  \leq 2\sqrt{2}.
\end{equation}
\end{theorem}

\begin{proof}[Proof sketch]
Classical local hidden variables can achieve at most $2$ (Bell's
bound).  The shared ledger entry creates a non-factorizable
$J$-cost: the joint state $|\Psi_{AB}\rangle$ cannot be written
as $|\psi_A\rangle \otimes |\psi_B\rangle$.  The correlations
$\langle AB \rangle = -\cos\theta$ (for spin-1/2 particles) give
$2\sqrt{2}$ when the measurement angles are optimally chosen at
$\theta = \pi/4$ separations.  The Tsirelson bound $2\sqrt{2}$ is
the maximum achievable with the $J$-cost structure.
\end{proof}

\begin{proposition}[No Signaling]
Despite the non-local correlations, no information can be
transmitted via entanglement.  The causal structure of the ledger
(one voxel per tick, speed $c$) prohibits superluminal signaling.
\end{proposition}

%% ============================================================
\section{Commutation Relations}
\label{sec:commutation}
%% ============================================================

\begin{theorem}[Canonical Commutation Relation]
\label{thm:commutation}
\begin{equation}
  \boxed{[\hat{x}, \hat{p}] = i\hbar,}
\end{equation}
where $\hbar = E_{\mathrm{coh}} \cdot \tau_0$.
\end{theorem}

\begin{proof}
Position $\hat{x}$ and momentum $\hat{p}$ correspond to
complementary recognition operators: $\hat{x}$ records the voxel
location of a recognition event; $\hat{p}$ records its frequency
mode.  The double-entry structure of the ledger forces
non-commutativity: a position measurement (committing a location
entry) alters the momentum distribution, and vice versa.

Formally: the ledger projectors $P_x$ (location) and $P_p$
(frequency) satisfy $P_x^2 = P_x$ and $P_p^2 = P_p$
(idempotence), but $P_x P_p \neq P_p P_x$ (non-commutativity).
The commutator $[P_x, P_p]$ has magnitude $\hbar = E_{\mathrm{coh}}
\cdot \tau_0$, the minimum action quantum on the ledger.  The
factor $i$ arises from the 8-tick complex phase structure
(Section~\ref{sec:hilbert}).
\end{proof}

\begin{corollary}[Uncertainty Principle]
\begin{equation}
  \Delta x \cdot \Delta p \geq \frac{\hbar}{2}.
\end{equation}
\end{corollary}

\begin{proof}
From the Robertson--Schr\"odinger inequality applied to
$[\hat{x}, \hat{p}] = i\hbar$:
$\Delta x \cdot \Delta p \geq \tfrac{1}{2}
|\langle [\hat{x}, \hat{p}] \rangle| = \hbar/2$.
\end{proof}

%% ============================================================
\section{The Quantum-to-Classical Transition}
\label{sec:classical}
%% ============================================================

\begin{theorem}[Decoherence from $J$-Cost Scaling]
\label{thm:decoherence}
For a system of $N$ particles, the decoherence time is:
\begin{equation}
  \tau_D \approx \frac{\tau_0\,E_{\mathrm{coh}}}{\Delta E}
  = \frac{\hbar}{\Delta E},
\end{equation}
where $\Delta E$ is the energy uncertainty.  For macroscopic
objects ($\Delta E \gg E_{\mathrm{coh}}$), $\tau_D \ll \tau_0$:
decoherence is effectively instantaneous.
\end{theorem}

\begin{proof}[Proof sketch]
The $J$-cost of an entangled $N$-particle state scales as $N^2$,
while a product state scales as $N$.  The cost differential
$\Delta J \sim N(N-1)$ drives decoherence at rate
$\Gamma_D \sim \Delta J / \tau_0$.  For $N \sim 10^{23}$:
$\Gamma_D \sim 10^{46}/\tau_0$, giving $\tau_D \sim 10^{-46}
\tau_0$---far below any observable timescale.
\end{proof}

\begin{remark}
There is no ``Heisenberg cut.''  The quantum-to-classical
transition is smooth, governed by the $N^2$ scaling.  Small
systems ($N \sim 1$) are quantum; large systems
($N \gg 1$) are classical; intermediate systems show partial
decoherence.
\end{remark}

%% ============================================================
\section{Comparison with Interpretations}
\label{sec:comparison}
%% ============================================================

\begin{center}
\renewcommand{\arraystretch}{1.3}
\begin{tabular}{lccccc}
\hline
\textbf{Feature} & \textbf{Copen.} & \textbf{MWI} &
\textbf{Bohm} & \textbf{GRW} & \textbf{RS} \\
\hline
Born rule & Postulated & Contested & Equivariant &
Derived & Derived \\
$\mathbb{C}$ & Assumed & Assumed & Assumed &
Assumed & Forced \\
Collapse & Observer & None & Epistemic &
Stochastic & Commit \\
Definite & Observer & All & Particles &
Random & Discrete \\
Free params & 0 & 0 & 0 & 2 & 0 \\
\hline
\end{tabular}
\end{center}

%% ============================================================
\section{Predictions and Falsifiers}
\label{sec:predictions}
%% ============================================================

\begin{enumerate}
  \item \textbf{Born rule exact}: $P = |A|^2$ at all scales.  Any
  violation falsifies the $J$-cost derivation.

  \item \textbf{No real QM}: Experiments requiring complex
  amplitudes (e.g., entanglement swapping across independent
  sources) should always confirm $\mathbb{C}$ over $\mathbb{R}$.

  \item \textbf{Tsirelson bound}: The maximum CHSH value is
  exactly $2\sqrt{2}$.  Exceeding this bound would falsify the
  ledger constraint model.

  \item \textbf{Decoherence time}: $\tau_D = \hbar/\Delta E$.
  Precision decoherence measurements should match this scaling.

  \item \textbf{No macroscopic superpositions}: Superpositions of
  $N \gtrsim 10^{10}$ particles should be unobservable.
\end{enumerate}

%% ============================================================
\section{Conclusion}
\label{sec:conclusion}
%% ============================================================

All five foundational features of quantum mechanics are derived
from $J(x) = \tfrac{1}{2}(x + x^{-1}) - 1$:
\begin{enumerate}[nosep]
  \item Born rule: unique probability function satisfying phase
  invariance, normalization, and continuity.
  \item $\mathbb{C}$: forced by the 8-tick half-period
  $e^{i\pi} = -1$.
  \item Measurement: ledger commit at epoch boundaries;
  no separate postulate.
  \item Entanglement: shared ledger constraint; Bell violation
  up to $2\sqrt{2}$.
  \item $[\hat{x}, \hat{p}] = i\hbar$: double-entry
  non-commutativity with minimum action $\hbar$.
\end{enumerate}

%% ============================================================
\begin{thebibliography}{99}

\bibitem{Bell1987}
J.~S.~Bell, \emph{Speakable and Unspeakable in Quantum Mechanics}
(Cambridge University Press, 1987).

\bibitem{Gleason}
A.~M.~Gleason, ``Measures on the Closed Subspaces of a Hilbert
Space,'' J.\ Math.\ Mech.\ \textbf{6}, 885 (1957).

\bibitem{RealQM2022}
M.-O.~Renou \emph{et al.}, ``Quantum Theory Based on Real
Numbers Can Be Experimentally Falsified,'' Nature \textbf{600},
625 (2021).

\bibitem{Zurek2003}
W.~H.~Zurek, ``Decoherence, Einselection, and the Quantum Origins
of the Classical,'' Rev.\ Mod.\ Phys.\ \textbf{75}, 715 (2003).

\end{thebibliography}

\end{document}
